% ============================================================
%  ACTIVIDAD: CONTROL DE CALIDAD EN EL AULA
%  Pinguinos Marinela. Fragmento independiente, para insertarlo donde convenga:
%
%      % ============================================================
%  ACTIVIDAD: CONTROL DE CALIDAD EN EL AULA
%  Pinguinos Marinela. Fragmento independiente, para insertarlo donde convenga:
%
%      % ============================================================
%  ACTIVIDAD: CONTROL DE CALIDAD EN EL AULA
%  Pinguinos Marinela. Fragmento independiente, para insertarlo donde convenga:
%
%      % ============================================================
%  ACTIVIDAD: CONTROL DE CALIDAD EN EL AULA
%  Pinguinos Marinela. Fragmento independiente, para insertarlo donde convenga:
%
%      \input{actividad-control-calidad}
%
%  Necesita el mismo preambulo de semana01.tex (usa \destaca y \Rc).
%
%  DONDE COLOCARLO. Las dos mitades de la actividad rinden en momentos distintos:
%    - la discusion del diseno (las dos primeras diapositivas) cae bien en la
%      unidad 5, produccion de datos;
%    - el analisis cae bien en la unidad 7, que es donde se aprende la prueba
%      contra un valor de referencia, que es exactamente la pregunta principal.
%  Se puede capturar en la unidad 5 y analizar en la 7, con los mismos datos.
% ============================================================
\begin{frame}{Actividad: un control de calidad en el aula}
  \begin{columns}[T]
    \begin{column}{0.55\textwidth}
      Cada quien recibe un \textbf{pingüino Marinela}. Antes de comerlo, hay que pesarlo.
      \vspace{0.3cm}

      A diferencia de un pan artesanal, este es un \textbf{producto industrial}: se
      supone que todas las piezas son iguales, y el empaque \textbf{declara un
      contenido neto}.
      \vspace{0.3cm}

      Ese número declarado es el \textbf{valor de referencia} de todo el análisis.
      No lo inventamos nosotros: viene impreso en la envoltura.
    \end{column}
    \begin{column}{0.42\textwidth}
      \destaca{\small La pregunta del control de calidad:
      \vspace{0.2cm}

      ¿El proceso \textbf{entrega lo que declara}? ¿Y con cuánta \textbf{variación} de
      pieza a pieza?}
      \vspace{0.3cm}

      {\footnotesize Existe una norma mexicana, la \textbf{NOM-002-SCFI}, que fija
      justamente esto: las tolerancias del contenido neto y los \textbf{planes de
      muestreo} para verificarlo. Es muestreo y prueba de hipótesis con fuerza de ley.}
    \end{column}
  \end{columns}
\end{frame}
% ============================================================
\begin{frame}{Lo que hay que acordar antes de medir}
  \centering
  {\scriptsize
  \begin{tabular}{@{}p{2.9cm}p{4.4cm}p{5.0cm}@{}}
    \toprule
    \textbf{Decisión} & \textbf{Qué hay que acordar} & \textbf{Por qué importa} \\
    \midrule
    \textbf{Definición operacional} & ¿El diámetro en su parte más ancha? ¿La altura
      con o sin la base de papel? &
      Sin esto cada quien mide otra cosa y los datos no son comparables \\
    \textbf{Instrumento y precisión} & Báscula de $0.1$ g para la masa, regla para el
      diámetro. Se anota la precisión &
      La precisión es parte del dato, y decide qué diferencias se pueden ver \\
    \textbf{Dos observadores} & Cada pieza la miden \emph{dos} personas distintas &
      Es la única forma de saber si la medida depende de quién mide \\
    \textbf{A ciegas} & El segundo no ve la lectura del primero &
      Si la ve, la copia, y el acuerdo entre ambos sería un artefacto \\
    \textbf{Orden contrabalanceado} & En la mitad de las piezas mide primero A; en la
      otra mitad, primero B &
      Al manipularla se desprende chocolate: la segunda masa sale más baja \\
    \textbf{Faltantes} & Anotar \texttt{NA} \textbf{y su causa} &
      Los faltantes tienen causa, y preguntarla es parte del análisis \\
    \bottomrule
  \end{tabular}}
  \vspace{0.35cm}

  \raggedright
  \destaca{\small Nótese qué se asigna al azar y qué no. El \textbf{peso de la pieza}
  ya venía dado por la fábrica: respecto a eso, esto es un estudio \textbf{observacional}.
  Lo que sí decidimos nosotros al azar es el \textbf{procedimiento de medición}, y
  sobre eso sí se puede hablar de causa.
  \vspace{0.15cm}

  Es la misma lógica de la cata de galletas (orden al azar, evaluación a ciegas),
  pero aquí la medición es \textbf{objetiva}: la hace una báscula, no un paladar.}
\end{frame}
% ============================================================
\begin{frame}{¿Qué preguntas se pueden responder?}
  \centering
  {\small
  \begin{tabular}{@{}p{7.2cm}p{4.8cm}@{}}
    \toprule
    \textbf{Pregunta} & \textbf{Técnica que la resuelve} \\
    \midrule
    ¿La masa promedio coincide con el contenido declarado? &
      Una muestra contra un valor de referencia {\scriptsize\color{grisT}(unidad 7)} \\
    ¿Cuánta variación hay de pieza a pieza? &
      Desviación estándar y coeficiente de variación {\scriptsize\color{grisT}(unidad 2)} \\
    ¿Qué proporción de piezas queda \emph{por debajo} de lo declarado? &
      Estimación de una proporción {\scriptsize\color{grisT}(unidad 9)} \\
    ¿Coinciden los dos observadores en la misma pieza? &
      Prueba $t$ pareada {\scriptsize\color{grisT}(unidad 8)} \\
    ¿Pierde masa al ser manipulada? &
      Prueba $t$ pareada {\scriptsize\color{grisT}(unidad 8)} \\
    ¿Difieren dos lotes distintos? &
      Dos muestras independientes {\scriptsize\color{grisT}(unidad 8)} \\
    ¿Cuántas piezas harían falta para detectar medio gramo? &
      Tamaño de muestra y potencia {\scriptsize\color{grisT}(unidad 7)} \\
    \midrule
    \textcolor{rojoAc}{¿Cuál báscula es \emph{más precisa}?} &
      \textcolor{rojoAc}{\textbf{No se puede con este diseño}} \\
    \bottomrule
  \end{tabular}}
  \vspace{0.3cm}

  \raggedright
  \destaca{\small La última merece atención: precisión es \textbf{variabilidad al
  repetir}, y aquí cada pieza se pesa una sola vez en cada báscula. Para contestarla
  habría que pesar la misma pieza varias veces con cada una. Es un \textbf{cambio de
  diseño}, no una técnica más sofisticada.}
\end{frame}
% ============================================================
\begin{frame}{Por qué se espera que pesen un poco de más}
  Antes de ver los datos vale la pena predecir el resultado. Hay una \textbf{asimetría
  económica} en el llenado:
  \vspace{0.35cm}

  \begin{columns}[T]
    \begin{column}{0.48\textwidth}
      \destaca{\textbf{Llenar de menos}\\[4pt]
      {\small Incumple lo declarado. Riesgo de sanción y de reclamo del consumidor.}}
    \end{column}
    \begin{column}{0.48\textwidth}
      \destaca{\textbf{Llenar de más}\\[4pt]
      {\small Cumple, pero \emph{cuesta}: es producto regalado, multiplicado por
      millones de piezas.}}
    \end{column}
  \end{columns}
  \vspace{0.4cm}

  \destaca{Por eso un proceso bien ajustado no apunta al valor declarado, sino
  \textbf{un poco por arriba}: lo suficiente para que casi ninguna pieza quede debajo,
  y no más. Eso convierte la pregunta en \textbf{unilateral}:
  $H_0: \mu = $ declarado contra $H_a: \mu > $ declarado.
  \vspace{0.2cm}

  Y cambia la interpretación. Si resulta significativo, la conclusión \textbf{no} es
  ``la empresa engaña'': es ``el proceso apunta por arriba de lo declarado, como era
  de esperarse''. Lo interesante para calidad no es tanto el promedio como la
  \textbf{variación}, porque es la variación la que produce piezas fuera de norma.}
\end{frame}

%
%  Necesita el mismo preambulo de semana01.tex (usa \destaca y \Rc).
%
%  DONDE COLOCARLO. Las dos mitades de la actividad rinden en momentos distintos:
%    - la discusion del diseno (las dos primeras diapositivas) cae bien en la
%      unidad 5, produccion de datos;
%    - el analisis cae bien en la unidad 7, que es donde se aprende la prueba
%      contra un valor de referencia, que es exactamente la pregunta principal.
%  Se puede capturar en la unidad 5 y analizar en la 7, con los mismos datos.
% ============================================================
\begin{frame}{Actividad: un control de calidad en el aula}
  \begin{columns}[T]
    \begin{column}{0.55\textwidth}
      Cada quien recibe un \textbf{pingüino Marinela}. Antes de comerlo, hay que pesarlo.
      \vspace{0.3cm}

      A diferencia de un pan artesanal, este es un \textbf{producto industrial}: se
      supone que todas las piezas son iguales, y el empaque \textbf{declara un
      contenido neto}.
      \vspace{0.3cm}

      Ese número declarado es el \textbf{valor de referencia} de todo el análisis.
      No lo inventamos nosotros: viene impreso en la envoltura.
    \end{column}
    \begin{column}{0.42\textwidth}
      \destaca{\small La pregunta del control de calidad:
      \vspace{0.2cm}

      ¿El proceso \textbf{entrega lo que declara}? ¿Y con cuánta \textbf{variación} de
      pieza a pieza?}
      \vspace{0.3cm}

      {\footnotesize Existe una norma mexicana, la \textbf{NOM-002-SCFI}, que fija
      justamente esto: las tolerancias del contenido neto y los \textbf{planes de
      muestreo} para verificarlo. Es muestreo y prueba de hipótesis con fuerza de ley.}
    \end{column}
  \end{columns}
\end{frame}
% ============================================================
\begin{frame}{Lo que hay que acordar antes de medir}
  \centering
  {\scriptsize
  \begin{tabular}{@{}p{2.9cm}p{4.4cm}p{5.0cm}@{}}
    \toprule
    \textbf{Decisión} & \textbf{Qué hay que acordar} & \textbf{Por qué importa} \\
    \midrule
    \textbf{Definición operacional} & ¿El diámetro en su parte más ancha? ¿La altura
      con o sin la base de papel? &
      Sin esto cada quien mide otra cosa y los datos no son comparables \\
    \textbf{Instrumento y precisión} & Báscula de $0.1$ g para la masa, regla para el
      diámetro. Se anota la precisión &
      La precisión es parte del dato, y decide qué diferencias se pueden ver \\
    \textbf{Dos observadores} & Cada pieza la miden \emph{dos} personas distintas &
      Es la única forma de saber si la medida depende de quién mide \\
    \textbf{A ciegas} & El segundo no ve la lectura del primero &
      Si la ve, la copia, y el acuerdo entre ambos sería un artefacto \\
    \textbf{Orden contrabalanceado} & En la mitad de las piezas mide primero A; en la
      otra mitad, primero B &
      Al manipularla se desprende chocolate: la segunda masa sale más baja \\
    \textbf{Faltantes} & Anotar \texttt{NA} \textbf{y su causa} &
      Los faltantes tienen causa, y preguntarla es parte del análisis \\
    \bottomrule
  \end{tabular}}
  \vspace{0.35cm}

  \raggedright
  \destaca{\small Nótese qué se asigna al azar y qué no. El \textbf{peso de la pieza}
  ya venía dado por la fábrica: respecto a eso, esto es un estudio \textbf{observacional}.
  Lo que sí decidimos nosotros al azar es el \textbf{procedimiento de medición}, y
  sobre eso sí se puede hablar de causa.
  \vspace{0.15cm}

  Es la misma lógica de la cata de galletas (orden al azar, evaluación a ciegas),
  pero aquí la medición es \textbf{objetiva}: la hace una báscula, no un paladar.}
\end{frame}
% ============================================================
\begin{frame}{¿Qué preguntas se pueden responder?}
  \centering
  {\small
  \begin{tabular}{@{}p{7.2cm}p{4.8cm}@{}}
    \toprule
    \textbf{Pregunta} & \textbf{Técnica que la resuelve} \\
    \midrule
    ¿La masa promedio coincide con el contenido declarado? &
      Una muestra contra un valor de referencia {\scriptsize\color{grisT}(unidad 7)} \\
    ¿Cuánta variación hay de pieza a pieza? &
      Desviación estándar y coeficiente de variación {\scriptsize\color{grisT}(unidad 2)} \\
    ¿Qué proporción de piezas queda \emph{por debajo} de lo declarado? &
      Estimación de una proporción {\scriptsize\color{grisT}(unidad 9)} \\
    ¿Coinciden los dos observadores en la misma pieza? &
      Prueba $t$ pareada {\scriptsize\color{grisT}(unidad 8)} \\
    ¿Pierde masa al ser manipulada? &
      Prueba $t$ pareada {\scriptsize\color{grisT}(unidad 8)} \\
    ¿Difieren dos lotes distintos? &
      Dos muestras independientes {\scriptsize\color{grisT}(unidad 8)} \\
    ¿Cuántas piezas harían falta para detectar medio gramo? &
      Tamaño de muestra y potencia {\scriptsize\color{grisT}(unidad 7)} \\
    \midrule
    \textcolor{rojoAc}{¿Cuál báscula es \emph{más precisa}?} &
      \textcolor{rojoAc}{\textbf{No se puede con este diseño}} \\
    \bottomrule
  \end{tabular}}
  \vspace{0.3cm}

  \raggedright
  \destaca{\small La última merece atención: precisión es \textbf{variabilidad al
  repetir}, y aquí cada pieza se pesa una sola vez en cada báscula. Para contestarla
  habría que pesar la misma pieza varias veces con cada una. Es un \textbf{cambio de
  diseño}, no una técnica más sofisticada.}
\end{frame}
% ============================================================
\begin{frame}{Por qué se espera que pesen un poco de más}
  Antes de ver los datos vale la pena predecir el resultado. Hay una \textbf{asimetría
  económica} en el llenado:
  \vspace{0.35cm}

  \begin{columns}[T]
    \begin{column}{0.48\textwidth}
      \destaca{\textbf{Llenar de menos}\\[4pt]
      {\small Incumple lo declarado. Riesgo de sanción y de reclamo del consumidor.}}
    \end{column}
    \begin{column}{0.48\textwidth}
      \destaca{\textbf{Llenar de más}\\[4pt]
      {\small Cumple, pero \emph{cuesta}: es producto regalado, multiplicado por
      millones de piezas.}}
    \end{column}
  \end{columns}
  \vspace{0.4cm}

  \destaca{Por eso un proceso bien ajustado no apunta al valor declarado, sino
  \textbf{un poco por arriba}: lo suficiente para que casi ninguna pieza quede debajo,
  y no más. Eso convierte la pregunta en \textbf{unilateral}:
  $H_0: \mu = $ declarado contra $H_a: \mu > $ declarado.
  \vspace{0.2cm}

  Y cambia la interpretación. Si resulta significativo, la conclusión \textbf{no} es
  ``la empresa engaña'': es ``el proceso apunta por arriba de lo declarado, como era
  de esperarse''. Lo interesante para calidad no es tanto el promedio como la
  \textbf{variación}, porque es la variación la que produce piezas fuera de norma.}
\end{frame}

%
%  Necesita el mismo preambulo de semana01.tex (usa \destaca y \Rc).
%
%  DONDE COLOCARLO. Las dos mitades de la actividad rinden en momentos distintos:
%    - la discusion del diseno (las dos primeras diapositivas) cae bien en la
%      unidad 5, produccion de datos;
%    - el analisis cae bien en la unidad 7, que es donde se aprende la prueba
%      contra un valor de referencia, que es exactamente la pregunta principal.
%  Se puede capturar en la unidad 5 y analizar en la 7, con los mismos datos.
% ============================================================
\begin{frame}{Actividad: un control de calidad en el aula}
  \begin{columns}[T]
    \begin{column}{0.55\textwidth}
      Cada quien recibe un \textbf{pingüino Marinela}. Antes de comerlo, hay que pesarlo.
      \vspace{0.3cm}

      A diferencia de un pan artesanal, este es un \textbf{producto industrial}: se
      supone que todas las piezas son iguales, y el empaque \textbf{declara un
      contenido neto}.
      \vspace{0.3cm}

      Ese número declarado es el \textbf{valor de referencia} de todo el análisis.
      No lo inventamos nosotros: viene impreso en la envoltura.
    \end{column}
    \begin{column}{0.42\textwidth}
      \destaca{\small La pregunta del control de calidad:
      \vspace{0.2cm}

      ¿El proceso \textbf{entrega lo que declara}? ¿Y con cuánta \textbf{variación} de
      pieza a pieza?}
      \vspace{0.3cm}

      {\footnotesize Existe una norma mexicana, la \textbf{NOM-002-SCFI}, que fija
      justamente esto: las tolerancias del contenido neto y los \textbf{planes de
      muestreo} para verificarlo. Es muestreo y prueba de hipótesis con fuerza de ley.}
    \end{column}
  \end{columns}
\end{frame}
% ============================================================
\begin{frame}{Lo que hay que acordar antes de medir}
  \centering
  {\scriptsize
  \begin{tabular}{@{}p{2.9cm}p{4.4cm}p{5.0cm}@{}}
    \toprule
    \textbf{Decisión} & \textbf{Qué hay que acordar} & \textbf{Por qué importa} \\
    \midrule
    \textbf{Definición operacional} & ¿El diámetro en su parte más ancha? ¿La altura
      con o sin la base de papel? &
      Sin esto cada quien mide otra cosa y los datos no son comparables \\
    \textbf{Instrumento y precisión} & Báscula de $0.1$ g para la masa, regla para el
      diámetro. Se anota la precisión &
      La precisión es parte del dato, y decide qué diferencias se pueden ver \\
    \textbf{Dos observadores} & Cada pieza la miden \emph{dos} personas distintas &
      Es la única forma de saber si la medida depende de quién mide \\
    \textbf{A ciegas} & El segundo no ve la lectura del primero &
      Si la ve, la copia, y el acuerdo entre ambos sería un artefacto \\
    \textbf{Orden contrabalanceado} & En la mitad de las piezas mide primero A; en la
      otra mitad, primero B &
      Al manipularla se desprende chocolate: la segunda masa sale más baja \\
    \textbf{Faltantes} & Anotar \texttt{NA} \textbf{y su causa} &
      Los faltantes tienen causa, y preguntarla es parte del análisis \\
    \bottomrule
  \end{tabular}}
  \vspace{0.35cm}

  \raggedright
  \destaca{\small Nótese qué se asigna al azar y qué no. El \textbf{peso de la pieza}
  ya venía dado por la fábrica: respecto a eso, esto es un estudio \textbf{observacional}.
  Lo que sí decidimos nosotros al azar es el \textbf{procedimiento de medición}, y
  sobre eso sí se puede hablar de causa.
  \vspace{0.15cm}

  Es la misma lógica de la cata de galletas (orden al azar, evaluación a ciegas),
  pero aquí la medición es \textbf{objetiva}: la hace una báscula, no un paladar.}
\end{frame}
% ============================================================
\begin{frame}{¿Qué preguntas se pueden responder?}
  \centering
  {\small
  \begin{tabular}{@{}p{7.2cm}p{4.8cm}@{}}
    \toprule
    \textbf{Pregunta} & \textbf{Técnica que la resuelve} \\
    \midrule
    ¿La masa promedio coincide con el contenido declarado? &
      Una muestra contra un valor de referencia {\scriptsize\color{grisT}(unidad 7)} \\
    ¿Cuánta variación hay de pieza a pieza? &
      Desviación estándar y coeficiente de variación {\scriptsize\color{grisT}(unidad 2)} \\
    ¿Qué proporción de piezas queda \emph{por debajo} de lo declarado? &
      Estimación de una proporción {\scriptsize\color{grisT}(unidad 9)} \\
    ¿Coinciden los dos observadores en la misma pieza? &
      Prueba $t$ pareada {\scriptsize\color{grisT}(unidad 8)} \\
    ¿Pierde masa al ser manipulada? &
      Prueba $t$ pareada {\scriptsize\color{grisT}(unidad 8)} \\
    ¿Difieren dos lotes distintos? &
      Dos muestras independientes {\scriptsize\color{grisT}(unidad 8)} \\
    ¿Cuántas piezas harían falta para detectar medio gramo? &
      Tamaño de muestra y potencia {\scriptsize\color{grisT}(unidad 7)} \\
    \midrule
    \textcolor{rojoAc}{¿Cuál báscula es \emph{más precisa}?} &
      \textcolor{rojoAc}{\textbf{No se puede con este diseño}} \\
    \bottomrule
  \end{tabular}}
  \vspace{0.3cm}

  \raggedright
  \destaca{\small La última merece atención: precisión es \textbf{variabilidad al
  repetir}, y aquí cada pieza se pesa una sola vez en cada báscula. Para contestarla
  habría que pesar la misma pieza varias veces con cada una. Es un \textbf{cambio de
  diseño}, no una técnica más sofisticada.}
\end{frame}
% ============================================================
\begin{frame}{Por qué se espera que pesen un poco de más}
  Antes de ver los datos vale la pena predecir el resultado. Hay una \textbf{asimetría
  económica} en el llenado:
  \vspace{0.35cm}

  \begin{columns}[T]
    \begin{column}{0.48\textwidth}
      \destaca{\textbf{Llenar de menos}\\[4pt]
      {\small Incumple lo declarado. Riesgo de sanción y de reclamo del consumidor.}}
    \end{column}
    \begin{column}{0.48\textwidth}
      \destaca{\textbf{Llenar de más}\\[4pt]
      {\small Cumple, pero \emph{cuesta}: es producto regalado, multiplicado por
      millones de piezas.}}
    \end{column}
  \end{columns}
  \vspace{0.4cm}

  \destaca{Por eso un proceso bien ajustado no apunta al valor declarado, sino
  \textbf{un poco por arriba}: lo suficiente para que casi ninguna pieza quede debajo,
  y no más. Eso convierte la pregunta en \textbf{unilateral}:
  $H_0: \mu = $ declarado contra $H_a: \mu > $ declarado.
  \vspace{0.2cm}

  Y cambia la interpretación. Si resulta significativo, la conclusión \textbf{no} es
  ``la empresa engaña'': es ``el proceso apunta por arriba de lo declarado, como era
  de esperarse''. Lo interesante para calidad no es tanto el promedio como la
  \textbf{variación}, porque es la variación la que produce piezas fuera de norma.}
\end{frame}

%
%  Necesita el mismo preambulo de semana01.tex (usa \destaca y \Rc).
%
%  DONDE COLOCARLO. Las dos mitades de la actividad rinden en momentos distintos:
%    - la discusion del diseno (las dos primeras diapositivas) cae bien en la
%      unidad 5, produccion de datos;
%    - el analisis cae bien en la unidad 7, que es donde se aprende la prueba
%      contra un valor de referencia, que es exactamente la pregunta principal.
%  Se puede capturar en la unidad 5 y analizar en la 7, con los mismos datos.
% ============================================================
\begin{frame}{Actividad: un control de calidad en el aula}
  \begin{columns}[T]
    \begin{column}{0.55\textwidth}
      Cada quien recibe un \textbf{pingüino Marinela}. Antes de comerlo, hay que pesarlo.
      \vspace{0.3cm}

      A diferencia de un pan artesanal, este es un \textbf{producto industrial}: se
      supone que todas las piezas son iguales, y el empaque \textbf{declara un
      contenido neto}.
      \vspace{0.3cm}

      Ese número declarado es el \textbf{valor de referencia} de todo el análisis.
      No lo inventamos nosotros: viene impreso en la envoltura.
    \end{column}
    \begin{column}{0.42\textwidth}
      \destaca{\small La pregunta del control de calidad:
      \vspace{0.2cm}

      ¿El proceso \textbf{entrega lo que declara}? ¿Y con cuánta \textbf{variación} de
      pieza a pieza?}
      \vspace{0.3cm}

      {\footnotesize Existe una norma mexicana, la \textbf{NOM-002-SCFI}, que fija
      justamente esto: las tolerancias del contenido neto y los \textbf{planes de
      muestreo} para verificarlo. Es muestreo y prueba de hipótesis con fuerza de ley.}
    \end{column}
  \end{columns}
\end{frame}
% ============================================================
\begin{frame}{Lo que hay que acordar antes de medir}
  \centering
  {\scriptsize
  \begin{tabular}{@{}p{2.9cm}p{4.4cm}p{5.0cm}@{}}
    \toprule
    \textbf{Decisión} & \textbf{Qué hay que acordar} & \textbf{Por qué importa} \\
    \midrule
    \textbf{Definición operacional} & ¿El diámetro en su parte más ancha? ¿La altura
      con o sin la base de papel? &
      Sin esto cada quien mide otra cosa y los datos no son comparables \\
    \textbf{Instrumento y precisión} & Báscula de $0.1$ g para la masa, regla para el
      diámetro. Se anota la precisión &
      La precisión es parte del dato, y decide qué diferencias se pueden ver \\
    \textbf{Dos observadores} & Cada pieza la miden \emph{dos} personas distintas &
      Es la única forma de saber si la medida depende de quién mide \\
    \textbf{A ciegas} & El segundo no ve la lectura del primero &
      Si la ve, la copia, y el acuerdo entre ambos sería un artefacto \\
    \textbf{Orden contrabalanceado} & En la mitad de las piezas mide primero A; en la
      otra mitad, primero B &
      Al manipularla se desprende chocolate: la segunda masa sale más baja \\
    \textbf{Faltantes} & Anotar \texttt{NA} \textbf{y su causa} &
      Los faltantes tienen causa, y preguntarla es parte del análisis \\
    \bottomrule
  \end{tabular}}
  \vspace{0.35cm}

  \raggedright
  \destaca{\small Nótese qué se asigna al azar y qué no. El \textbf{peso de la pieza}
  ya venía dado por la fábrica: respecto a eso, esto es un estudio \textbf{observacional}.
  Lo que sí decidimos nosotros al azar es el \textbf{procedimiento de medición}, y
  sobre eso sí se puede hablar de causa.
  \vspace{0.15cm}

  Es la misma lógica de la cata de galletas (orden al azar, evaluación a ciegas),
  pero aquí la medición es \textbf{objetiva}: la hace una báscula, no un paladar.}
\end{frame}
% ============================================================
\begin{frame}{¿Qué preguntas se pueden responder?}
  \centering
  {\small
  \begin{tabular}{@{}p{7.2cm}p{4.8cm}@{}}
    \toprule
    \textbf{Pregunta} & \textbf{Técnica que la resuelve} \\
    \midrule
    ¿La masa promedio coincide con el contenido declarado? &
      Una muestra contra un valor de referencia {\scriptsize\color{grisT}(unidad 7)} \\
    ¿Cuánta variación hay de pieza a pieza? &
      Desviación estándar y coeficiente de variación {\scriptsize\color{grisT}(unidad 2)} \\
    ¿Qué proporción de piezas queda \emph{por debajo} de lo declarado? &
      Estimación de una proporción {\scriptsize\color{grisT}(unidad 9)} \\
    ¿Coinciden los dos observadores en la misma pieza? &
      Prueba $t$ pareada {\scriptsize\color{grisT}(unidad 8)} \\
    ¿Pierde masa al ser manipulada? &
      Prueba $t$ pareada {\scriptsize\color{grisT}(unidad 8)} \\
    ¿Difieren dos lotes distintos? &
      Dos muestras independientes {\scriptsize\color{grisT}(unidad 8)} \\
    ¿Cuántas piezas harían falta para detectar medio gramo? &
      Tamaño de muestra y potencia {\scriptsize\color{grisT}(unidad 7)} \\
    \midrule
    \textcolor{rojoAc}{¿Cuál báscula es \emph{más precisa}?} &
      \textcolor{rojoAc}{\textbf{No se puede con este diseño}} \\
    \bottomrule
  \end{tabular}}
  \vspace{0.3cm}

  \raggedright
  \destaca{\small La última merece atención: precisión es \textbf{variabilidad al
  repetir}, y aquí cada pieza se pesa una sola vez en cada báscula. Para contestarla
  habría que pesar la misma pieza varias veces con cada una. Es un \textbf{cambio de
  diseño}, no una técnica más sofisticada.}
\end{frame}
% ============================================================
\begin{frame}{Por qué se espera que pesen un poco de más}
  Antes de ver los datos vale la pena predecir el resultado. Hay una \textbf{asimetría
  económica} en el llenado:
  \vspace{0.35cm}

  \begin{columns}[T]
    \begin{column}{0.48\textwidth}
      \destaca{\textbf{Llenar de menos}\\[4pt]
      {\small Incumple lo declarado. Riesgo de sanción y de reclamo del consumidor.}}
    \end{column}
    \begin{column}{0.48\textwidth}
      \destaca{\textbf{Llenar de más}\\[4pt]
      {\small Cumple, pero \emph{cuesta}: es producto regalado, multiplicado por
      millones de piezas.}}
    \end{column}
  \end{columns}
  \vspace{0.4cm}

  \destaca{Por eso un proceso bien ajustado no apunta al valor declarado, sino
  \textbf{un poco por arriba}: lo suficiente para que casi ninguna pieza quede debajo,
  y no más. Eso convierte la pregunta en \textbf{unilateral}:
  $H_0: \mu = $ declarado contra $H_a: \mu > $ declarado.
  \vspace{0.2cm}

  Y cambia la interpretación. Si resulta significativo, la conclusión \textbf{no} es
  ``la empresa engaña'': es ``el proceso apunta por arriba de lo declarado, como era
  de esperarse''. Lo interesante para calidad no es tanto el promedio como la
  \textbf{variación}, porque es la variación la que produce piezas fuera de norma.}
\end{frame}
